% TOP_PROBLEM: {"id": 3000033, "title": "Equitable list colouring", "category_id": 2, "category": {"id": 2, "name": "combinatorics", "display_name": "Combinatorics", "description": "Counting problems, graph theory, discrete structures.", "slug": "combinatorics", "order_index": 2, "created_at": "2026-07-31T15:26:25.671Z"}, "status": "open", "problem_number": "AMR-029-0033", "set_id": 13}
% FIRST_POSTED: 2026-10-01
% ENTRY_KIND: statement_only
% UnsolvedMath ID 3000033; source catalogue code AMR-029-0033
% !TeX program = lualatex
% Definition and sources only; this file is not an LLM solution attempt.
\documentclass[11pt,a4paper]{article}
\usepackage[margin=25mm]{geometry}
\usepackage{fontspec}
\setmainfont{lmroman10-regular.otf}[BoldFont=lmroman10-bold.otf,
 ItalicFont=lmroman10-italic.otf,BoldItalicFont=lmroman10-bolditalic.otf]
\usepackage{amsmath,amssymb,mathtools}
\usepackage{unicode-math}
\setmathfont{latinmodern-math.otf}
\AtBeginDocument{\let\setminus\smallsetminus}
\usepackage{xurl,hyperref}
\hypersetup{colorlinks=true,urlcolor=blue}
\setcounter{secnumdepth}{0}
\setlength{\parindent}{0pt}
\setlength{\parskip}{5pt}
\setlength{\emergencystretch}{3em}
\begin{document}
\section*{Equitable list colouring}\label{3000033}
{\small UnsolvedMath ID 3000033 · AMR-029-0033 · Combinatorics · AMR Open Problem Lists}
\subsection{Definitions and mathematical statement}
Let $G=(V,E)$ be a finite simple graph, $n=|V|\ge1$, and
$\Delta(G)=\max_{v\in V}d_G(v)$. A $k$-list assignment supplies a set
$L(v)$ of exactly $k$ distinct colour names to each vertex. Lists
need not agree, and their union may contain more than $k$ colours.
A proper list colouring chooses $f(v)\in L(v)$ so that adjacent
vertices receive different colours.

Such a colouring is equitable in the list-colouring sense if
\[
                       |\{v:f(v)=a\}|\le\left\lceil\frac nk\right\rceil
                                      \quad\text{for every colour }a.
\]
The conjecture asserts that for every integer $k\ge\Delta(G)+1$ and
every $k$-list assignment, a proper colouring with this bound exists.

Only the upper size bound is imposed on each colour class. A matching
lower bound is inappropriate for arbitrary lists, since some colour
names may occur at only one vertex. In particular, the assertion does
not require the use of a common palette of exactly $k$ colours.
The lists are chosen adversarially and independently at the vertices;
the conclusion must hold for each assignment. Ordinary equitable
colourability, where every list is the same, is thus only a special
case of the problem.
\subsection{Short English statement}
Can every graph be properly coloured from arbitrary lists of
size at least one more than its maximum degree, without any colour being used too often?
\subsection{Sources}
\begin{itemize}
\item[S1] ULAM.AI and the UnsolvedMath Contributors, supplied problems.json, record 3000033 (AMR-029-0033), AMR Open Problem Lists. Catalogue identity and source transcription; CC BY 4.0.
\url{https://huggingface.co/datasets/ulamai/UnsolvedMath}.
\item[S2] Egres Open - List of open problems, Open-problem page: Equitable list colouring. Original source indexed by AMR-029-0033.
\url{https://oldlemon.cs.elte.hu/egres/open/List_of_open_problems}.
\item[S3] EGRES Open, Equitable list colouring. Individual source page and explanatory remarks.
\url{https://lemon.cs.elte.hu/egres/open/Equitable_list_colouring}.
\end{itemize}
\end{document}
